\begin{appendices}
\chapter{Грамматика поддерживаемого языка B}

Ниже приведена укрупненная грамматика поддерживаемого ядра языка B. Запись не повторяет дословно историческую нотацию руководства, а адаптирует ее к реализации компилятора.

Листинг А.1 --- Укрупненная грамматика внешних определений.
\begin{verbatim}
program     ::= { definition }

definition  ::= name [ "[" [ constant ] "]" ] [ initializer-list ] ";"
              | name "(" [ parameter-list ] ")" compound-statement

initializer-list ::= initializer { "," initializer }
initializer      ::= constant | name
parameter-list   ::= parameter { "," parameter }
parameter        ::= name
\end{verbatim}

Листинг А.2 --- Грамматика операторов.
\begin{verbatim}
statement ::= "auto" auto-list ";"
            | "extrn" name-list ";"
            | name ":" statement
            | "case" constant ":" statement
            | "{" { statement } "}"
            | "if" "(" expression ")" statement [ "else" statement ]
            | "while" "(" expression ")" statement
            | "switch" expression statement
            | "goto" expression ";"
            | "return" [ "(" expression ")" | expression ] ";"
            | [ expression ] ";"

auto-list ::= auto-decl { "," auto-decl }
auto-decl ::= name [ constant | "[" [ constant ] "]" ]
name-list ::= name { "," name }
\end{verbatim}

Листинг А.3 --- Грамматика выражений.
\begin{verbatim}
expression ::= assignment
assignment ::= conditional [ assign-op assignment ]
conditional ::= binary-expression
              | binary-expression "?" expression ":" conditional

primary    ::= name
             | constant
             | string
             | "(" expression ")"

postfix    ::= primary { "(" [ argument-list ] ")"
                       | "[" expression "]"
                       | "++"
                       | "--" }

prefix     ::= postfix
             | "&" prefix
             | "*" prefix
             | "-" prefix
             | "!" prefix
             | "~" prefix
             | "++" prefix
             | "--" prefix

binary-expression ::= prefix { binary-op prefix }
\end{verbatim}

При синтаксическом анализе выражений используется таблица приоритетов операторов. Присваивание является правоассоциативным. Тернарный оператор также связывается справа налево. Остальные бинарные операторы обрабатываются с учетом их приоритета и левой ассоциативности.
\end{appendices}
